%% ==================== 页面与排版 ====================
% XeLaTeX提供原生Unicode支持和fontspec字体加载能力，是中文论文排版的基础引擎要求
\RequirePackage{iftex}
\ifxetex\else
  \ifluatex
    \ClassWarning{cauc_thesis}{LuaLaTeX detected. This template is optimized for XeLaTeX. Some features may not work correctly with LuaLaTeX.}%
  \else
    \ClassError{cauc_thesis}{%
      This template requires XeLaTeX (recommended) or LuaLaTeX.%
      \MessageBreak
      Current engine is not compatible. Please switch your compiler.%
      \MessageBreak
      For TeX Live: latexmk -xelatex yourfile.tex%
    }{%
      Switch to XeLaTeX or LuaLaTeX engine to compile this template.%
    }%
    \stop
  \fi
\fi
\RequirePackage{float}
\RequirePackage{microtype}
\ifxetex
  \microtypecontext{spacing}
\fi
\ifluatex
  \microtypecontext{spacing}
\fi
\RequirePackage{setspace}
\RequirePackage{afterpage}
\RequirePackage[top=\cauc@topmargin,bottom=\cauc@bottommargin,left=\cauc@leftmargin,right=\cauc@rightmargin]{geometry}

%% ==================== 图形与表格 ====================
\RequirePackage{graphicx}
\RequirePackage{xcolor}
\graphicspath{{figures/}}
\RequirePackage{tabularx}
\RequirePackage{makecell}
\RequirePackage{array}
\RequirePackage{caption}
\RequirePackage{booktabs}
\RequirePackage{longtable}
\RequirePackage{multirow}
\RequirePackage{subcaption}
\RequirePackage{rotating}

%% ==================== 字体与文字 ====================
\RequirePackage{fontspec}
\RequirePackage{ragged2e}
\RequirePackage{ulem}
\RequirePackage{zhnumber}

%% ==================== 数学与算法 ====================
\RequirePackage{bm, amsmath, amssymb}
\def\cauc@@temp{times}%
\ifx\cauc@mathfont\cauc@@temp
  \RequirePackage{newtxmath}
\fi
\RequirePackage{siunitx}
\sisetup{group-separator={\,}, group-minimum-digits=4}
\ifcauc@graduate
\RequirePackage[ruled,linesnumbered]{algorithm2e}
\SetAlgorithmName{算法}{算法}{算法列表}
\fi

%% ==================== 代码与列表 ====================
\RequirePackage{enumitem}
\RequirePackage{listings}
\definecolor{cauc@keywordcolor}{RGB}{0,0,180}
\definecolor{cauc@commentcolor}{RGB}{0,128,0}
\definecolor{cauc@stringcolor}{RGB}{180,0,0}
\definecolor{cauc@green}{RGB}{28,172,0}
\definecolor{cauc@lilas}{RGB}{170,55,241}
\lstset{
	basicstyle=\zihao{-4}\ttfamily,
	keywordstyle=\color{cauc@keywordcolor}\bfseries,
	commentstyle=\color{cauc@commentcolor},
	stringstyle=\color{cauc@stringcolor},
	numbers=left,
	numberstyle=\tiny\color{gray},
	stepnumber=1,
	numbersep=8pt,
	backgroundcolor=\color{white},
	frame=single,
	rulecolor=\color{black!30},
	breaklines=true,
	breakatwhitespace=false,
	tabsize=2,
	showstringspaces=false,
	captionpos=b,
	extendedchars=true,
	texcl=true,
	escapeinside={(*@}{@*)}
}

%% ==================== 文档结构 ====================
\RequirePackage{xparse}
\RequirePackage{calc}
\RequirePackage{titlesec}
\RequirePackage{fancyhdr}
\RequirePackage{appendix}
\RequirePackage{titletoc}

%% ==================== TikZ绘图 ====================
\RequirePackage{tikz}

%% ==================== 引用与交叉 ====================
%% 引用宏包加载顺序：natbib -> hyperref -> cleveref（不可更改）
% natbib须最先加载，hyperref须在natbib之后、cleveref之前加载，cleveref须最后加载
% 新增宏包请在此段之前添加，或使用\AtBeginDocument避免干扰加载顺序
\ifcauc@bachelor
\RequirePackage[super,square,compress,comma]{natbib}
\else
\RequirePackage[sort&compress]{natbib}
\fi
\RequirePackage[colorlinks,
	anchorcolor=black,
	linkcolor=black,
	citecolor=black,
	urlcolor=black,
	unicode,
	pdfencoding=auto
]{hyperref}
\hypersetup{
	pdftitle={\cauc@Title},
	pdfauthor={\cauc@StudentName},
	pdfsubject={\ifcauc@bachelor Bachelor\else\ifcauc@doctor Doctor\else Master\fi\fi\ Thesis, \cauc@Department},
	pdfkeywords={\cauc@Keywords},
	pdfcreator={cauc\_thesis template with XeLaTeX}
}
\RequirePackage{bookmark}
% 消除hyperref PDF书签中的间距命令警告（~、\enspace、\quad等在PDF字符串中不支持）
\pdfstringdefDisableCommands{%
  \def~{ }%
  \def\enspace{ }%
  \def\quad{ }%
}
\IfFileExists{xurl.sty}{\RequirePackage{xurl}}{}
%% cleveref须在hyperref之后加载
\RequirePackage{cleveref}

%% ==================== 条件编译 ====================
% PDF/A模式需加载pdfx宏包以生成符合存档标准的PDF文件
\ifcauc@pdfa
  \RequirePackage[a-2b,mathxmp]{pdfx}
\fi

% 研究生论文需要双语标题(bicaption)、符号说明(nomencl)和术语表(glossaries)，本科不需要
\ifcauc@graduate
  \RequirePackage{bicaption}
  \captionsetup[figure][bi-second]{name=Fig.}
  \captionsetup[table][bi-second]{name=Tab.}
  \captionsetup{labelsep=space}
  \captionsetup[table]{labelsep=space}
  \RequirePackage{nomencl}
  \RequirePackage{pifont}
  \RequirePackage[toc,acronym,nonumberlist]{glossaries}
\fi


